\section{参数方程*}
\subsection{参数方程的概念}
\clearpage
\subsection{圆锥曲线的参数方程}
\begin{theorem}[圆的参数方程]
    以 \( O(x_0, y_0) \)为圆心，\(r\)为 半径的 \( \bigodot O  \) 的参数方程为
    \[\left\{
    \begin{gathered}
        x = x_0 + r \cos \alpha, \\
        y = y_0 + r \sin \alpha,
    \end{gathered}\right.
    \ (\alpha\text{为参数})
    \]
\end{theorem}


\begin{example}
    已知 \(x^2 + y^2 - 2y - 3 = 0\)，求 \(2x^2 + 2y^2 - 3x\) 的取值范围。
\end{example}

\begin{example}
    求\(f(x)=\dfrac{\sin x-3}{\cos x -1}\)的值域。
\end{example}

\begin{example}
    已知\(P\)是\(\bigodot C:(x-1)^2+y^2=1\)上一动点，两定点\(M(-1,0), N(0,-1)\)，求\(|PM|^2+|PN|^2\)的取值范围。
\end{example}
\begin{exercise}
    已知\(P\)是\(\bigodot C:x^2+y^2=4\)上一动点，两定点\(A(0,-4), B(2,0)\)，求\(\dfrac{|PB|}{|PA|}\)的取值范围。
\end{exercise}
\begin{theorem}[椭圆的参数方程]
标准方程为\(\dfrac{x^2}{a^2}+\dfrac{y^2}{b^2}=1\)的椭圆的参数方程为
    \[\left\{
    \begin{gathered}
        x = a \cos \alpha, \\
        y = b \sin \alpha,
    \end{gathered}\right.
    \ (\alpha\text{为参数})
    \]
\end{theorem}
\begin{corollary}[椭圆参数方程中参数的几何意义]
    
\end{corollary}
\begin{example}
    已知椭圆 \(\dfrac{x^2}{25} + \dfrac{y^2}{9} = 1\)，直线 \(l: 4x - 5y + 40 = 0\)。椭圆上是否存在一点，它到直线 \(l\) 的距离最小？最小距离是多少？
\end{example}


\begin{exercise}
    求\(f(x)=\dfrac{\sin x-\cos x}{2\cos x -1}\)的值域。
\end{exercise}

\begin{example}
    设椭圆 \(E: \dfrac{x^2}{a^2} + \dfrac{y^2}{b^2} = 1 (a > b > 0)\) 的左、右焦点为 \(F_1, F_2\)，右顶点为 \(A\)，已知点 \(P\) 在椭圆 \(E\) 上，若 \(\angle F_1PF_2 = 90^\circ\)，\(\angle PAF_2 = 45^\circ\)，求椭圆 \(E\) 的离心率。
\end{example}
\begin{example}
    
\end{example}
其他圆锥曲线自然也有参数方程。
\begin{theorem}[双曲线的参数方程]
标准方程为\(\dfrac{x^2}{a^2}-\dfrac{y^2}{b^2}=1\)的双曲线的参数方程为
    \[\left\{
    \begin{gathered}
        x = a \sec \alpha, \\
        y = b \tan \alpha,
    \end{gathered}\right.
    \ (\alpha\text{为参数})
    \]
\end{theorem}

\[\left\{
    \begin{gathered}
        x = a \sinh \alpha, \\
        y = b \cosh \alpha,
    \end{gathered}\right.
    \ (\alpha\text{为参数})
    \]

\begin{theorem}
标准方程为\(y^2=2px\)的抛物线的参数方程为
\[\left\{
    \begin{gathered}
        x = 2pt^2, \\
        y = 2pt,
    \end{gathered}\right.
    \ (t\text{为参数})
    \]
\end{theorem}


\clearpage



\subsection{直线的参数方程}
\begin{theorem}[直线的参数方程]
经过点 \( M_0(x_0, y_0) \)，倾斜角为 \( \alpha \) 的直线 \( l \) 的参数方程为
\[\left\{
\begin{gathered}
    x = x_0 + t \cos \alpha, \\
    y = y_0 + t \sin \alpha,
\end{gathered}\right.
\ (t\text{为参数})
\]
\end{theorem}
\begin{corollary}[直线的参数方程中\(t\)的几何意义]
    
\end{corollary}
\begin{example}
    完成直线的一般方程和参数方程的互化。
    
    \begin{enumerate}
        \item \(4x+3y-5=0\)
        \item \(y=x\)
        \item \(y=5x+2\)
    \end{enumerate}
\end{example}
\clearpage

